% =========================================================
% Ordered Sets: Partial, Strict, and Total Orders
% =========================================================

\section{Ordered Sets}
\label{sec:ordered-sets}

\begin{remark*}[Toolkit: Ordered Sets]
Order is a relation, not a property.  A set carries no intrinsic ordering;
an order is an additional structure imposed on the set by specifying which
pairs of elements stand in the ordering relation.  The same underlying set
can support many distinct orders simultaneously.  The atomic notions
developed here are:
\begin{itemize}
  \item the \textbf{partial order} --- the minimal framework in which order
        makes sense;
  \item the \textbf{strict partial order} --- the asymmetric companion to the
        partial order, and the native language of open-ball conditions and
        limits;
  \item the \textbf{total order} --- the specialisation in which every pair of
        elements is comparable.
\end{itemize}
\end{remark*}

% ---------------------------------------------------------
% def:lattice-order-partial-order
% ---------------------------------------------------------

\begin{definitionbox}{Definition (Partial Order)}
\begin{definition}[Partial Order]
\label{def:lattice-order-partial-order}
A relation $\leq$ on a set $S$ is a \emph{partial order} if it satisfies,
for all $a, b, c \in S$:
\begin{enumerate}
  \item \textbf{Reflexivity:} $a \leq a$.
  \item \textbf{Anti-symmetry:} $a \leq b \land b \leq a \Longrightarrow a = b$.
  \item \textbf{Transitivity:} $a \leq b \land b \leq c \Longrightarrow a \leq c$.
\end{enumerate}
The pair $(S, \leq)$ is called a \emph{partially ordered set}, or \emph{poset}.
\end{definition}
\end{definitionbox}

\begin{remark*}[Standard quantified statement]
\[
\operatorname{PartialOrder}(S, \leq)
\;\coloneqq\;
\bigl(\forall a \in S,\; a \leq a\bigr)
\;\land\;
\bigl(\forall a,b \in S,\; a \leq b \land b \leq a \Rightarrow a = b\bigr)
\;\land\;
\bigl(\forall a,b,c \in S,\; a \leq b \land b \leq c \Rightarrow a \leq c\bigr).
\]
\end{remark*}

\begin{remark*}[Negated quantified statement]
\[
\neg\operatorname{PartialOrder}(S, \leq)
\;\coloneqq\;
\bigl(\exists a \in S,\; a \not\leq a\bigr)
\;\lor\;
\bigl(\exists a,b \in S,\; a \leq b \land b \leq a \land a \neq b\bigr)
\;\lor\;
\bigl(\exists a,b,c \in S,\; a \leq b \land b \leq c \land a \not\leq c\bigr).
\]
\end{remark*}

\begin{remark*}[Interpretation]
Order is not an intrinsic quality of a set's elements --- it is an external
binary relation.  The three axioms together enforce that the relation be a
consistent, cycle-free, directed comparison.  Reflexivity says every element
is at least as large as itself.  Anti-symmetry is the formal engine behind
squeeze arguments: if $a \leq b$ and $b \leq a$ then $a$ and $b$ cannot
be distinct.  Transitivity ensures that the ordering chains: comparisons
can be composed without breaking consistency.  The same set of objects can
carry many distinct partial orders simultaneously --- standard, reverse,
or lexicographic --- each defining a different relational geometry.
\end{remark*}

\begin{dependencies}
\begin{itemize}
  \item \hyperref[def:algebraic-partial-order]{Partial Order}
\end{itemize}
\end{dependencies}
% ---------------------------------------------------------
% def:lattice-order-strict-partial-order
% ---------------------------------------------------------

\begin{definitionbox}{Definition}
\begin{definition}[Strict Partial Order]
\label{def:lattice-order-strict-partial-order}
A relation $<$ on a set $S$ is a \emph{strict partial order} if it satisfies,
for all $a, b, c \in S$:
\begin{enumerate}
  \item \textbf{Irreflexivity:} $\neg(a < a)$.
  \item \textbf{Asymmetry:} $a < b \Longrightarrow \neg(b < a)$.
  \item \textbf{Transitivity:} $a < b \land b < c \Longrightarrow a < c$.
\end{enumerate}
\end{definition}
\end{definitionbox}

\begin{remark*}[Standard quantified statement]
\[
\operatorname{StrictPartialOrder}(S, {<})
\;\coloneqq\;
\bigl(\forall a \in S,\; \neg(a < a)\bigr)
\;\land\;
\bigl(\forall a,b \in S,\; a < b \Rightarrow \neg(b < a)\bigr)
\;\land\;
\bigl(\forall a,b,c \in S,\; a < b \land b < c \Rightarrow a < c\bigr).
\]
\end{remark*}

\begin{remark*}[Interpretation]
The strict order is the companion of the non-strict order, eliminating
reflexivity and replacing anti-symmetry with the stronger asymmetry.
Every strict partial order determines a partial order by
$a \leq b \Longleftrightarrow (a < b \lor a = b)$, and every partial order
determines a strict partial order by $a < b \Longleftrightarrow (a \leq b
\land a \neq b)$.  In metric spaces and analysis, strict orders are the
natural language of open conditions: open balls, open intervals, and limit
definitions all use $<$ rather than $\leq$ to exclude boundary contact.
\end{remark*}

\begin{dependencies}
\begin{itemize}
  \item \hyperref[def:lattice-order-partial-order]{Partial Order}
\end{itemize}
\end{dependencies}

% ---------------------------------------------------------
% def:lattice-order-total-order
% ---------------------------------------------------------

\begin{definitionbox}{Definition (Total Order)}
\begin{definition}[Total Order]
\label{def:lattice-order-total-order}
A \hyperref[def:lattice-order-partial-order]{partial order} $\leq$ on a set $S$ is a
\emph{total order} (or \emph{linear order}) if it additionally satisfies
the \emph{law of trichotomy}: for all $a, b \in S$, exactly one of
\[
a < b, \qquad a = b, \qquad b < a
\]
holds.  The pair $(S, \leq)$ is called a \emph{totally ordered set}
or a \emph{chain}.
\end{definition}
\end{definitionbox}

\begin{remark*}[Standard quantified statement]
\[
\operatorname{TotalOrder}(S, \leq)
\;\coloneqq\;
\operatorname{PartialOrder}(S, \leq)
\;\land\;
\forall a, b \in S,\;
(a < b) \oplus (a = b) \oplus (b < a),
\]
where $\oplus$ denotes exclusive or.  Equivalently, comparability suffices:
$\forall a,b \in S,\; a \leq b \lor b \leq a$.
\end{remark*}

\begin{remark*}[Interpretation]
A partial order permits elements to be incomparable: it is possible that
neither $a \leq b$ nor $b \leq a$.  A total order removes this
possibility entirely --- every pair of elements is forced to be comparable.
The rational numbers $\mathbb{Q}$ and real numbers $\mathbb{R}$ under their
standard orderings are totally ordered fields.  By contrast, the power
set $\mathcal{P}(X)$ of any set $X$ with more than one element, ordered
by inclusion, is typically not a total order: the singletons $\{a\}$ and
$\{b\}$ with $a \neq b$ are incomparable.
\end{remark*}

\begin{dependencies}
\hyperref[def:lattice-order-partial-order]{Partial order},
\hyperref[def:lattice-order-strict-partial-order]{strict partial order}.
\end{dependencies}
